

\section{Wissenschaftliche Fragestellung} \label{cha:wissenschaftlicheFragestellung}
%\section{Research Question} \label{cha:researchQuestion}

Die in Kapitel \ref{cha:standDesWissens} durchgeführte Literaturanalyse verdeutlicht im Wesentlichen zwei Aspekte: (1) Es gibt eine Vielzahl an Studien und Paper, die sich mit LLMs und Programmieren allgemein oder im engeren Sinne mit dem Einsatz von LLMs zum Programmierenlernen beschäftigen. So fokussieren z.B. \textcite{RobotsAreComming} oder \textcite{KazemitabaarSupportLearners2023} Performance-Metriken wie Zeit, Korrektheit und Code-Qualität. Der Aspekt der User Experience bei der Interaktion mit einem LLM-Chatbot wird jedoch wenig bis selten betrachtet. Die Arbeiten von \textcite{CodeHelp}, \textcite{KazemitabaarCodeAid} und \textcite{KazemitabaarSupportLearners2023} nehmen hierbei eine Sonderstellung ein, da sie Einblicke in die Benutzerfreundlichkeit ihrer Tools bieten. Bei ersterem wurde eine Online-Umfrage durchgeführt, bei zweiterem Interviews sowie wöchentliche Umfragen und bei letzterem wurde Feedback mittels Likert-Skala-Fragen eingeholt. Ein standardisiertes Instrument wie der UEQ kam nicht zum Einsatz. Die Literaturanalyse verdeutlicht (2) auch, dass in Arbeiten, die den UEQ eingesetzen, es primär um die Evaluierung von LLM-Tools geht ohne spezifischen Bezug zur Programmierausbildung.

Daraus ergibt sich folgende Forschungslücke: LLM-Chatbots werden zwar als Hilfe für Programmiereinsteiger bereitgestellt, es fehlt jedoch eine standardisierte Auswertung der User Experience in diesem Bereich. Um diesem Umstand entgegenzuwirken, wird in der vorliegenden Arbeit ein LLM-Chatbot herangezogen -- in diesem Fall ChatGPT -- mit dessen Unterstützung das Programmieren-lernen bei Einsteigern gefördert werden soll. Dieser Prozess wird nicht nur aus didaktischer Sicht betrachtet werden, sondern auch aus dem Blickwinkel der User Experience.

Die folgenden Forschungsfragen werden im Laufe der Arbeit beantwortet:

\textbf{RQ1:} Inwieweit beeinflusst die Nutzung eines LLMs (ChatGPT) die Leistung von ProgrammieranfängerInnen bei der Ausarbeitung einer Prüfungsaufgabe im Vergleich zu einer Bearbeitung ohne LLM?

\textbf{RQ2:} Wie unterscheiden sich User Experience und subjektive Wahrnehmung der Prüfungssituation zwischen den beiden Versuchsgruppen, gemessen mit dem UEQ sowie offenen Reflexionsfragen?

Gemessen wird RQ1 anhand der Faktoren: Bearbeitungszeit, Anzahl der Programmstarts bis zur Abgabe und Korrektheit der Gesamtlösung. Die Bewertung der Korrektheit erfordert einen detaillierten und einheitlichen Bewertungsraster, anhand dessen der im Versuch produzierte Programmcode evaluiert werden kann.